\chapter{IBVS-MPC: o alvo na imagem dentro do custo}
\label{cap:ibvs_mpc}

\section{Motivação}
Nas leis dos capítulos anteriores, a centralização do alvo na imagem fica \emph{fora} do
controlador. O MPC (\autoref{cap:mpc}) calcula só a translação; o yaw
(\code{yaw\_mode: face\_target}) e o canal vertical (\code{vertical\_mode},
\autoref{cap:centralizacao}) são políticas geométricas aplicadas depois. Isso tem três limitações:
\begin{enumerate}
  \item o MPC não sabe que o yaw precisa acompanhar o alvo: uma manobra ótima em posição pode
        tirar o alvo do campo de visão enquanto o yaw ainda converge;
  \item com a câmera fixa no corpo, \emph{acelerar inclina a imagem} (a arfagem segue a direção
        do empuxo). A lei de translação não enxerga esse acoplamento;
  \item o yaw de \code{face\_target} é um apontamento estático, sem predição do alvo nem
        \emph{feedforward} de taxa, e fica sempre atrasado em relação a um alvo que gira.
\end{enumerate}
O \code{ibvs\_mpc} (\code{control/ibvs\_mpc\_controller.cpp}) põe a posição do alvo na imagem
no próprio problema de controle ótimo. É um MPC não linear \emph{perception-aware}: as coordenadas
de imagem são as \emph{features} do IBVS, e a profundidade (disponível na OakD-Lite) fixa a
distância.

Há duas formulações, escolhidas por \code{ibvs.formulation}. Este capítulo descreve primeiro a
\code{perception\_aware} (padrão), que segue o modelo \emph{global} de \cite{allibert2010vpc} e o
PAMPC de \cite{falanga2018pampc}. A \code{classic} (\autoref{sec:ibvs_classic}) é a câmera
virtual da literatura de quadrirrotores \cite{sheng2019msc,zhang2021rnmpc} e serve de linha de
base.

\section{Modelo de predição}
O estado é $x=[\vect p;\vect v;\psi]$ (posição relativa à posição atual, velocidade ENU, yaw ENU)
e a entrada é $u=[\vect a;r]$ (aceleração e taxa de yaw). A grade é a mesma do LTV-MPC,
$\Delta_k=\Ts g^k$ ($N=16$, $g=1{,}12$, $T_H\approx\SI{2,1}{s}$):
\begin{equation}
  \vect p_{k+1}=\vect p_k+\vect v_k\Delta_k+\tfrac12\vect a_k\Delta_k^2,\qquad
  \vect v_{k+1}=\vect v_k+\vect a_k\Delta_k,
\end{equation}
e o yaw tem a malha de taxa da PX4 como 1ª ordem com atraso ($r$ é a demanda de taxa):
\begin{equation}
  \dot\omega=\frac{r(t-\tau_d)-\omega}{\tau_\psi},\qquad\dot\psi=\omega,
  \qquad
  \begin{aligned}
    \omega_{k+1}&=\alpha_k\omega_k+(1-\alpha_k)r_k,\\
    \psi_{k+1}&=\psi_k+r_k\Delta_k+(\omega_k-r_k)\,\tau_\psi(1-\alpha_k),
  \end{aligned}
  \quad \alpha_k=e^{-\Delta_k/\tau_\psi}.
\end{equation}
O alvo é previsto como nas demais leis, com velocidade constante ou Singer
(\code{targetDisplacement}): $\vect t_k=\vect p_{rel}+\Delta\vect p_{alvo}(t_k)$.
O yaw inicial é o \emph{heading} medido (\code{VehicleLocalPosition.heading}) e $\omega_0$ é a taxa
medida (odometria). O atraso é compensado como num preditor de Smith: $(\psi_0,\omega_0)$ são
propagados pelos comandos já enviados nos últimos $\tau_d$ segundos.

\subsection{Identificação da malha de yaw (SITL)}
\label{sec:ibvs_yaw_id}
A primeira versão supunha $\dot\psi=r$ (integrador ideal). No SITL (\code{s3\_figure8}) o azimute
do alvo oscilava $\pm\SI{25}{\degree}$: a taxa comandada alternava entre $\pm\SI{1,5}{rad/s}$ com
período de \SI{1,2}{s}, um ciclo-limite.

Reprodução: \code{ros2\_ws/identify\_px4\_lag.py --yaw <sessões>}. A demanda de taxa é refeita como
a PX4 a forma, $u=\sat(2{,}8\,e_\psi+\dot\psi_{ff})$ (\code{MC\_YAW\_P} $=2{,}8$,
\code{MC\_YAWRATE\_MAX} $=\SI{200}{\degree/s}$, padrões da PX4 v1.18), e a taxa medida é a taxa de
Euler da odometria. O ajuste é de \emph{erro de saída}: 1ª ordem com atraso simulada no registro
inteiro, busca em grade de $\tau_\psi$ e $\tau_d$. O ajuste de um passo à frente (equação de erro)
foi descartado: no mesmo voo dava $\tau$ de 56 a \SI{320}{ms} conforme o passo de reamostragem, por
causa do ruído da derivada do \emph{heading}.

\begin{center}
\begin{tabular}{lccc}
  \toprule
  Sessões (SITL, 15/11) & $\tau_\psi$ & $\tau_d$ & resíduo / desvio da taxa \\
  \midrule
  ciclo-limite (figure-8 e zigzag, 1ª versão) & 0,28 s & 0,16 s & 0,17 / \SI{1,0}{rad/s} \\
  \code{ibvs\_mpc} com o modelo de yaw (6 voos) & 0,15--0,24 s & 0,06--0,14 s & 0,02 / \SI{0,16}{}--\SI{0,32}{rad/s} \\
  LTV + \code{face\_target} (8 voos) & 0,26--0,34 s & 0,02--0,10 s & 0,01--0,03 / \SI{0,15}{}--\SI{0,29}{rad/s} \\
  \bottomrule
\end{tabular}
\end{center}
Com pouca excitação, $\tau_\psi$ e $\tau_d$ trocam entre si e só a soma (0,29 a \SI{0,44}{s}) é bem
determinada. Os valores do modelo, $\tau_\psi=\SI{0,30}{s}$ e $\tau_d=\SI{0,16}{s}$, vêm dos voos de
ciclo-limite, os únicos com excitação para separar os dois; os resultados variam menos de
\SI{10}{ms} entre passos de 10 e \SI{50}{ms}. Com essa planta no teste unitário, o integrador ideal
reproduz a falha (\SI{12,8}{\degree} RMS de azimute, alvo parado). Só a 1ª ordem reduz para
\SI{3,9}{\degree}; com a compensação do atraso o erro vai a zero com alvo parado e a \SI{2,2}{\degree}
com o alvo em curva. Em hardware: degraus de taxa de guinada e o mesmo script.

\subsection{Câmera}
A atitude do corpo vem da planicidade diferencial: o eixo $z_B$ segue o empuxo e o yaw fixa o
eixo $x$,
\begin{equation}
  z_B=\frac{\vect a+g\,\vect e_z}{\|\vect a+g\,\vect e_z\|},\quad
  y_B=\frac{z_B\times x_\psi}{\|z_B\times x_\psi\|},\quad x_B=y_B\times z_B,\quad
  x_\psi=[\cos\psi,\ \sin\psi,\ 0]\T .
\end{equation}
Com \code{camera\_model: level} (gimbal em \code{stabilize}) a inclinação é ignorada,
$R_{WB}=R_z(\psi)$. A montagem é uma arfagem $\theta_c$ (\code{camera\_pitch}, $+$ para cima, lida
da TF estática como no \code{center\_image}): $R_{WC}=R_{WB}R_y(-\theta_c)$. As \emph{features} são
os ângulos esféricos do alvo na câmera (convenção do \code{TargetBearing}):
\begin{equation}
  \vect c = R_{WC}\T(\vect t_k-\vect p_k),\qquad
  s_{az}=\operatorname{atan2}(c_y,c_x),\qquad
  s_{el}=\operatorname{atan2}\!\left(c_z,\sqrt{c_x^2+c_y^2}\right).
\end{equation}
Ao contrário das coordenadas normalizadas $c_y/c_x$, os ângulos são suaves para qualquer direção
(inclusive com o alvo atrás), o que mantém o problema bem posto na aquisição.

\subsection{Refinamentos do modelo da câmera (comutáveis)}
\label{sec:ibvs_refinamentos}
Três correções entram na projeção. Cada uma tem uma chave \code{ibvs.enable.*} (aba Controle da
GCS); desligada, o modelo volta ao anterior.
\begin{itemize}
  \item \textbf{Posição da câmera} (\code{enable.camera\_offset}). A estimativa $\vect p_{rel}$ parte da
        origem do corpo, mas a câmera do x500\_depth fica em $(0{,}12;\ 0{,}03;\ 0{,}24)$~m (TF estática
        base $\to$ \code{camera\_link}). A \SI{5}{m}, os \SI{0,24}{m} de altura são \SI{2,8}{\degree} de
        elevação. A projeção passa a usar $\vect t_k-\vect p_k-R_{WB}\,\vect o_c$.
  \item \textbf{Atraso da inclinação} (\code{enable.tilt\_lag}, \code{tilt\_time\_constant}). A atitude
        não acompanha a aceleração planejada na hora: $\dot{\bar{\vect a}}=(\vect a-\bar{\vect a})/\tau_t$,
        e o corpo segue $\bar{\vect a}$. O estado inicial $\bar{\vect a}_0$ vem da atitude medida. O ajuste
        nos bags (\code{identify\_px4\_lag.py --tilt}, erro de saída) deu $\tau_t+\tau_d\approx\SI{0,18}{s}$
        no eixo $x$ (resíduo \SI{0,12}{m/s^2} contra \SI{0,18}{m/s^2} sem atraso, no zigzag). No eixo
        $y$ o modelo explica pouco (resíduo de 0,56 a \SI{0,95}{m/s^2}); causa não investigada.
  \item \textbf{Viés de inclinação} (\code{enable.tilt\_bias}, \code{tilt\_bias\_filter}). Em cruzeiro o
        drone fica inclinado para vencer o arrasto, com aceleração nula. O estimador
        $\vect d\leftarrow\vect d+\tfrac{\Ts}{\tau_b+\Ts}\,(\bar{\vect a}_{med}-\bar{\vect a}_{mod}-\vect d)$
        compara a inclinação medida com a explicada só pelos comandos ($\bar{\vect a}_{mod}$, nunca
        reiniciada pela medida) e mantém $\vect d$ constante no horizonte: o corpo segue
        $\bar{\vect a}+\vect d$. No teste unitário com arrasto (\SI{2,6}{\degree} de inclinação a
        \SI{1,5}{m/s}) a elevação RMS cai de \SI{1,98}{\degree} para \SI{0,10}{\degree}. No SITL, porém,
        piorou o figure-8 (\autoref{sec:ibvs_sitl}) e por isso vem \textbf{desligado por padrão}.
\end{itemize}
As demais chaves desligam partes já descritas: \code{yaw\_lag}, \code{yaw\_delay},
\code{los\_rotation\_reference}, \code{closing\_speed}, \code{fov\_constraint}, \code{image\_huber} e
\code{acquisition}.

\subsection{Gimbal no loop (\code{ibvs.gimbal\_in\_loop})}
\label{sec:ibvs_gimbal}
Com a câmera no gimbal (\autoref{sec:gimbal}), a arfagem e a guinada da \emph{câmera} viram graus de
liberdade de verdade. No corpo de um quadrotor não é assim: arfagem e rolagem são a direção do
empuxo, e o hover exige corpo nivelado. O MPC ganha dois estados, a direção do eixo óptico no mundo
$\vect c=[c_{az},c_{el}]$, e duas entradas, o \emph{setpoint} $\vect u_g$ enviado ao gimbal:
\begin{equation}
  \dot{\vect c}=\frac{\vect u_g-\vect c}{\tau_g},\qquad \tau_g=\frac{1}{k_{gimbal}}=\SI{0,17}{s},
  \qquad \|\vect u_g-\vect c\|_\infty\le\omega_{max}\tau_g,
\end{equation}
o seguidor do \code{gimbal\_controller} (ganho \code{inertia.gain}, taxa máxima
\code{inertia.max\_rate\_deg\_s}). A câmera é estabilizada, então a projeção usa
$R_{WC}=R_z(c_{az})R_y(-c_{el})$ e não depende da inclinação do corpo. O custo ganha
\begin{equation}
  w_{gy}(c_{az}-\psi)^2+w_{gp}(c_{el}-\theta_0)^2+w_{gl}\,[\text{limites das juntas}]_+^2
  +w_{gu}\|\vect u_g-\vect c\|^2 .
\end{equation}
Os dois primeiros termos dividem o trabalho: o gimbal corrige rápido, e o corpo (yaw e altitude)
corrige devagar até a junta voltar ao centro, longe dos limites. O limite de taxa é imposto por
projeção sequencial sobre os estados planejados. Sem ele, o primeiro \emph{setpoint} saltava mais de
$\pi$ e o seguidor, que usa o erro embrulhado, girava pelo lado errado (teste
\code{GimbalSetpointTurnsTheShortWayWithinTheRateLimit}).

O nó publica $\vect u_g$ em \code{/gimbal/pointing\_setpoint}, e o \code{gimbal\_controller} ganhou o
modo \code{mpc}, que aponta para esse \emph{setpoint} (sem o \emph{feedforward} de taxa, porque o MPC
já antecipa) e cai para \code{track} se ele ficar velho. A direção medida vem de
\code{/gimbal/state} com a atitude.

\section{Custo}
Por estágio ($k=1..N$, o último multiplicado por \code{terminal\_weight}):
\begin{multline}
  \ell_k = w_{img}\,\rho_\delta(s_{az}-s_{az}^*)^2+w_{img}\,\rho_\delta(s_{el}-s_{el}^*)^2
  + w_r\big(\|\vect t_k-\vect p_k\|-d\big)^2 + w_v\|\vect v_k-\vect v_{des,k}\|^2\\
  + w_{fov}\textstyle\sum_i\max\!\big(0,|s_i|-m\,\tfrac{\mathrm{FOV}_i}{2}\big)^2
  + \varrho\max(0,\|\vect v_k\|-v_{max})^2 + w_a\|\vect a_{k-1}\|^2+w_\psi r_{k-1}^2 .
\end{multline}
\begin{itemize}
  \item \textbf{Imagem} com pseudo-Huber $\rho_\delta$ (\code{image\_huber} $\delta=\SI{0,25}{rad}$):
        quadrático perto do centro e linear para erros grandes. Sem isso, com o alvo a $90^\circ$ o
        ótimo usava a \emph{inclinação} do corpo (aceleração de \SI{3}{m/s^2}) para girar a câmera,
        porque o yaw é limitado em taxa. O peso cai a $w_{img}\cdot$\code{acquisition\_weight} quando o
        alvo está fora da imagem agora: quem o traz de volta é o yaw.
  \item \textbf{Distância} $\|\vect t_k-\vect p_k\|-d$. Não há ponto cartesiano de referência
        (IBVS puro, \code{weight\_pos}$=0$): distância e imagem fixam a geometria relativa.
  \item \textbf{Velocidade}: velocidade do ponto de \emph{standoff} quando só o alvo gira a LOS
        (a mesma referência do LTV-MPC), no plano horizontal, mais uma velocidade de aproximação:
        \begin{equation}
          \vect v_{des}=\vect v_t-\tfrac{d}{\max(r,d)}(I-\vect l_h\vect l_h\T)\vect v_{t,h}
          +\sat_{v_c}\!\big(k_c(r-d)\big)\vect l .
        \end{equation}
        Copiar $\vect v_t$ inteiro brigava com a distância nas curvas (erro de \SI{0,3}{m} no teste de
        malha fechada). Remover também a componente vertical deixava de acompanhar alvos subindo,
        porque na vertical só a imagem centraliza. O termo $k_c(r-d)$ (\code{closing\_gain}) evita
        que o termo de velocidade freie a aproximação quando o drone está longe.
  \item \textbf{FOV} (\code{weight\_fov}, margem \code{fov\_margin}) é uma restrição suave
        ativa só com o alvo dentro da imagem. Ela mantém o alvo visível e não persegue um alvo já perdido.
\end{itemize}

\section{Solução numérica}
O problema é um OCP do \textbf{acados} (\code{aero\_ibvs[\_lag][\_gim]} em
\code{acados/aero\_ocps.py}), escrito em CasADi: estado
$[\vect p,\vect v,(\vect v_r),\bar{\vect a},\psi,\omega_\psi,(\vect c)]$, entradas
$[\vect a,r,(\vect u_g)]$, os blocos de 1ª ordem como EDOs e o custo como mínimos quadrados não
lineares (os resíduos acima, com os pesos como parâmetros de execução). \emph{Multiple shooting}
com integrador RK4, SQP com Hessiana de Gauss--Newton e busca linear por mérito, QP pelo HPIPM; as
derivadas são exatas (diferenciação automática do CasADi). Restrições:
$\|\vect a_k\|\le a_{max}$ e a taxa do seguidor do gimbal duras, $|r_k|\le r_{max}$ por caixa, e
FOV, $\|\vect v_k\|\le v_{max}$ e limites das juntas do gimbal brandas (folgas com penalidade
quadrática, pesos \code{weight\_fov}, \code{velocity\_penalty}, \code{weight\_gimbal\_limits}). O
\emph{warm start} usa a solução anterior deslocada, com no máximo \code{max\_iterations} = 6
iterações por ciclo. Antes de resolver, o yaw medido (e o azimute da câmera) é desembrulhado
para perto do plano anterior: um salto de $2\pi$ entre os estágios 0 e 1 levaria o Gauss--Newton a
girar o drone pelo lado longo.

\paragraph{Saída.} $\vect v=\vect v_{ref}+\vect a_0\Ts$ e $\vect a_0$ como nas demais leis, e
$\psi_{sp}=\psi_{med}+\Delta\psi_0$ (o incremento previsto do 1º estágio, sobre o \emph{heading}
\emph{medido}, para o termo P de yaw da PX4 ver só o passo deste ciclo) com \code{yawspeed}$=-r_0$ (NED) como \emph{feedforward} no
\code{TrajectorySetpoint}. Com esta lei, \code{yaw\_mode} e \code{vertical\_mode} são ignorados
(valem só em APPROACH, sem profundidade).

\section{Resultados em malha fechada (simulação do teste unitário)}
\code{test/test\_ibvs\_mpc.cpp} simula uma planta tipo PX4: malha de velocidade de 1ª ordem
($\tau=\SI{0,2}{s}$, \autoref{cap:latencia}), malha de yaw identificada (\autoref{sec:ibvs_yaw_id}) e câmera
no corpo com a atitude real (planicidade sobre a aceleração \emph{real}). O alvo faz uma curva a
\SI{1,5}{m/s} ($\omega=\SI{0,25}{rad/s}$) com subida senoidal ($\pm$\SI{0,4}{m/s}). As métricas são
tomadas após \SI{20}{s}:

\begin{center}
\begin{tabular}{lcccc}
  \toprule
  Lei & az RMS & el RMS & máx. fora do eixo & erro de distância RMS \\
  \midrule
  LTV-MPC + \code{face\_target} + \code{los}          & $5{,}15^\circ$ & $8{,}14^\circ$ & $12{,}6^\circ$ & \SI{0,04}{m} \\
  LTV-MPC + \code{face\_target} + \code{center\_image} & $5{,}21^\circ$ & $0{,}58^\circ$ & $5{,}5^\circ$ & \SI{0,04}{m} \\
  \code{ibvs\_mpc} (câmera no corpo)                   & $2{,}17^\circ$ & $0{,}98^\circ$ & $2{,}8^\circ$ & \SI{0,10}{m} \\
  \code{ibvs\_mpc}, câmera estabilizada (\code{level})  & $2{,}12^\circ$ & $0{,}79^\circ$ & $2{,}4^\circ$ & \SI{0,04}{m} \\
  \code{ibvs\_mpc}, gimbal no loop                     & $0{,}04^\circ$ & $0{,}05^\circ$ & $0{,}1^\circ$ & \SI{0,04}{m} \\
  \midrule
  \multicolumn{5}{l}{Câmera inclinada $-15^\circ$:}\\
  LTV-MPC + \code{face\_target} + \code{los}          & $5{,}36^\circ$ & $18{,}3^\circ$ & $27{,}0^\circ$ (sai do FOV) & \SI{0,04}{m} \\
  LTV-MPC + \code{face\_target} + \code{center\_image} & $5{,}51^\circ$ & $0{,}61^\circ$ & $5{,}8^\circ$ & \SI{0,04}{m} \\
  \code{ibvs\_mpc}                                     & $2{,}10^\circ$ & $1{,}23^\circ$ & $3{,}0^\circ$ & \SI{0,08}{m} \\
  \bottomrule
\end{tabular}
\end{center}

O ganho principal é horizontal. O yaw previsto, com \emph{feedforward} de taxa, reduz o atraso
de \SI{5}{\degree} do \code{face\_target} a \SI{2}{\degree}. Na vertical o \code{center\_image} continua um pouco
melhor, e o \code{ibvs\_mpc} paga com cerca de \SI{6}{cm} a mais de erro de distância: é a troca
entre imagem e distância nos pesos. Com o gimbal no loop o erro de imagem praticamente some, porque
a câmera deixa de depender do yaw e da inclinação do corpo. O custo computacional médio é de
\SIrange{1,6}{2,7}{ms} por ciclo (\SIrange{4}{6}{ms} com o gimbal; \SI{0,2}{ms} no LTV), folgado
para o período de \SI{50}{ms}.

\section{Resultados no SITL (PX4 + Gazebo)}
\label{sec:ibvs_sitl}
Cenários \code{s3\_figure8} e \code{s4\_zigzag}, \SI{45}{s}, 3 repetições (média; o desvio entre
repetições ficou abaixo de \SI{0,008}{rad}). O erro de imagem é o \emph{verdadeiro}: posição do centro
visível do alvo na câmera real, pelas TFs de verdade do Gazebo gravadas no bag
(\code{scripts/image\_error.py}); o RMS total inclui a aproximação inicial, e azimute/elevação são
tomados após \SI{15}{s}. Estes estudos (\code{ibvs\_mpc\_full}, \code{ibvs\_mpc\_vision\_full},
\code{ibvs\_mpc\_gimbal}) rodaram antes da correção da posição da câmera.

\begin{center}
\begin{tabular}{llcccc}
  \toprule
  Cenário & Lei & erro de imagem & az / el RMS & \code{rmse\_position} & $\int\|\vect a\|^2$ \\
  \midrule
  \multicolumn{6}{l}{Câmera fixa, detecção omnisciente}\\
  figure-8 & LTV + \code{face\_target} + \code{los}          & \SI{0,242}{rad} & $9{,}2^\circ$ / $9{,}6^\circ$ & \SI{6,64}{m} & 41 \\
           & LTV + \code{face\_target} + \code{center\_image} & \SI{0,179}{rad} & $9{,}0^\circ$ / $5{,}1^\circ$ & \SI{6,98}{m} & 80 \\
           & \code{ibvs\_mpc}                                 & \SI{0,106}{rad} & $3{,}1^\circ$ / $4{,}3^\circ$ & \SI{5,71}{m} & 102 \\
  zigzag   & LTV + \code{face\_target} + \code{los}          & \SI{0,099}{rad} & $4{,}5^\circ$ / $3{,}8^\circ$ & \SI{2,62}{m} & 18 \\
           & LTV + \code{face\_target} + \code{center\_image} & \SI{0,095}{rad} & $4{,}3^\circ$ / $4{,}3^\circ$ & \SI{2,44}{m} & 30 \\
           & \code{ibvs\_mpc}                                 & \SI{0,070}{rad} & $1{,}4^\circ$ / $3{,}0^\circ$ & \SI{1,42}{m} & 50 \\
  \midrule
  \multicolumn{6}{l}{Câmera fixa, visão (tracker ViT + profundidade)}\\
  figure-8 & LTV + \code{face\_target} + \code{center\_image} & \SI{0,195}{rad} & $10{,}1^\circ$ / $7{,}1^\circ$ & \SI{2,89}{m} & 57 \\
           & \code{ibvs\_mpc}                                 & \SI{0,134}{rad} & $6{,}0^\circ$ / $5{,}5^\circ$ & \SI{2,60}{m} & 60 \\
  zigzag   & LTV + \code{face\_target} + \code{center\_image} & \SI{0,139}{rad} & $6{,}7^\circ$ / $6{,}2^\circ$ & \SI{2,82}{m} & 41 \\
           & \code{ibvs\_mpc}                                 & \SI{0,116}{rad} & $4{,}5^\circ$ / $4{,}9^\circ$ & \SI{1,44}{m} & 59 \\
  \midrule
  \multicolumn{6}{l}{Gimbal, detecção omnisciente}\\
  figure-8 & LTV + \code{center\_image}, gimbal \code{stabilize}        & \SI{0,156}{rad} & $9{,}4^\circ$ / $1{,}1^\circ$ & & \\
           & \code{ibvs\_mpc} \code{level}, gimbal \code{stabilize}     & \SI{0,050}{rad} & $2{,}8^\circ$ / $1{,}4^\circ$ & & \\
           & \code{ibvs\_mpc} com gimbal no loop (\code{mpc})           & \SI{0,026}{rad} & $1{,}1^\circ$ / $1{,}1^\circ$ & \SI{5,64}{m} & 102 \\
  zigzag   & LTV + \code{center\_image}, gimbal \code{stabilize}        & \SI{0,069}{rad} & $4{,}5^\circ$ / $0{,}9^\circ$ & & \\
           & \code{ibvs\_mpc} \code{level}, gimbal \code{stabilize}     & \SI{0,030}{rad} & $1{,}4^\circ$ / $1{,}0^\circ$ & & \\
           & \code{ibvs\_mpc} com gimbal no loop (\code{mpc})           & \SI{0,016}{rad} & $0{,}4^\circ$ / $0{,}9^\circ$ & & \\
  \bottomrule
\end{tabular}
\end{center}
O \code{ibvs\_mpc} reduz o erro de imagem em todos os casos, com rastreamento acima de
\SI{99,7}{\percent} e nenhuma perda de alvo na visão. O custo é o esforço de controle (até
\SI{65}{\percent} acima do \code{center\_image} no zigzag). O gimbal no loop é o melhor resultado, mas
\textbf{falta a linha de base} LTV + gimbal em \code{track} (o gimbal centraliza sozinho) e o teste
com visão: \code{ibvs\_mpc\_gimbal\_loop.yaml} e \code{ibvs\_mpc\_gimbal\_loop\_vision.yaml} estão
prontos e não foram rodados.

\paragraph{A/B dos refinamentos (câmera fixa, omnisciente).}
\begin{center}
\begin{tabular}{llccc}
  \toprule
  Chave & Cenário & desligada & ligada & elevação média \\
  \midrule
  \code{tilt\_lag}       & figure-8 & \SI{0,116}{rad} & \SI{0,115}{rad} & \\
                         & zigzag   & \SI{0,078}{rad} & \SI{0,076}{rad} & \\
  \code{camera\_offset}  & figure-8 & \SI{0,108}{rad} & \SI{0,102}{rad} & $-2{,}7^\circ\to+0{,}5^\circ$ \\
                         & zigzag   & \SI{0,072}{rad} & \SI{0,072}{rad} & $-1{,}5^\circ\to+1{,}4^\circ$ \\
  \code{tilt\_bias}      & figure-8 & \SI{0,100}{rad} & \SI{0,118}{rad} & $+0{,}4^\circ\to+1{,}4^\circ$ \\
  \bottomrule
\end{tabular}
\end{center}
O atraso da inclinação é neutro (fica ligado por ser inofensivo). A posição da câmera remove o viés
de elevação do figure-8. O estimador do viés de inclinação \emph{piora} o figure-8, apesar do bom
resultado no teste unitário: com o alvo sempre em curva, o resíduo que ele filtra não é arrasto
constante. Fica desligado por padrão; o zigzag (onde sobra $+1{,}4^\circ$ de viés) não foi concluído.
As duas primeiras linhas usaram a medida sem o centro visível do alvo, por isso o nível difere.

\section{Formulação clássica (\code{ibvs.formulation: classic})}
\label{sec:ibvs_classic}
\code{control/classic\_ibvs\_mpc\_controller.cpp}. Reproduz \cite{sheng2019msc} e
\cite{zhang2021rnmpc}: câmera virtual $\mathcal V$ (a real sem rolagem e arfagem; mantém a
inclinação de montagem), eixos ópticos, \emph{features} no plano virtual, Euler com o período de
controle, custo $\sum(\vect x-\vect x_{ref})\T Q(\vect x-\vect x_{ref})+\sum\vect u\T R\vect u$, entradas
$\vect u=[\vect f;\dot\psi]$ com $\|\vect f\|\le a_{max}$, $|\dot\psi|\le r_{max}$, e rumo
$\psi_{sp}=\psi_m+\dot\psi T_s$ (Eq.~3.24 de Sheng).
\begin{itemize}
  \item \code{features: moments}: $\vect s=\vect P/Z^*$, $\dot{\vect s}=-\vect\omega\times\vect s-(\vect v-\vect v_t)/Z^*$
        (Eq.~2.28 de Sheng para câmera frontal); \code{points}: $[X/Z,\ Y/Z,\ Z/Z^*]$ pela matriz de
        interação de \cite{chaumette2006}. $s_3$ vem da profundidade (igual a $\sqrt{a^*/a}$ para alvo
        de tamanho constante).
  \item Velocidade no referencial girante: $\dot{\vect v}=\vect f-\vect\omega\times\vect v$, com
        $\vect\omega=\dot\psi\,R_{W\mathcal V}\T\vect e_z$.
  \item \code{target\_motion} (adaptação, ligada): $\vect v-\vect v_t$ na cinemática e $\vect v_t$ como
        referência de velocidade (Sheng: alvo parado, $\vect v_{ref}=0$).
  \item Restrição de imagem de \cite{allibert2010vpc} e \cite{chan_ibvs_mpc} como penalidade
        (\code{weight\_image\_bounds}); nas referências é dura (SQP).
  \item Não modela: malhas da PX4 (yaw e velocidade ideais), a inclinação na imagem real, a posição
        da câmera, o gimbal. \code{mpc.lag} não se aplica.
\end{itemize}

\paragraph{Sintonia por varredura.} \code{ClassicIbvsMpc.DISABLED\_TuningSweep} (192 combinações de
$N_p\in\{10,20\}$, $N_c\in\{1,3,N_p\}$, $R_{\dot\psi}\in\{0{,}05;0{,}5;2;8\}$, $R_f\in\{0{,}05;0{,}5\}$,
$r_{max}\in\{1{,}5;0{,}4\}$) no teste de malha fechada abaixo. Conclusões: $R_{\dot\psi}=0{,}05$ põe o
rumo em ciclo-limite ($\approx\SI{12}{\degree}$ RMS de azimute) contra a malha de yaw da PX4, e
$0{,}5$ o elimina; $N_c=1$ (Sheng) quadruplica o erro de distância; $N_c=3$ iguala $N_c=N_p$ com
um décimo do tempo; $N_p=20$ ($\SI{1}{s}$) reduz o erro de distância de 0,68 para \SI{0,28}{m}
em relação a 10; $r_{max}$ não muda nada (não satura). $Q$ é o de Sheng.

\paragraph{Comparação no teste de malha fechada} (\code{ClassicIbvsMpc.ClosedLoopComparison},
mesma planta e alvo da \autoref{sec:ibvs_sitl}):
\begin{center}
\begin{tabular}{lccccc}
  \toprule
  & az RMS & el RMS & máx. fora do eixo & distância RMS & tempo/ciclo \\
  \midrule
  perception-aware & $2{,}17^\circ$ & $0{,}98^\circ$ & $2{,}8^\circ$ & \SI{0,10}{m} & \SI{2,1}{ms} \\
  clássica, momentos & $0{,}12^\circ$ & $1{,}34^\circ$ & $3{,}5^\circ$ & \SI{0,28}{m} & \SI{0,4}{ms} \\
  clássica, ponto & $0{,}12^\circ$ & $1{,}48^\circ$ & $4{,}1^\circ$ & \SI{0,28}{m} & \SI{0,4}{ms} \\
  clássica, alvo parado & $2{,}70^\circ$ & $6{,}22^\circ$ & $9{,}4^\circ$ & \SI{0,96}{m} & \SI{0,6}{ms} \\
  clássica, $R_{\dot\psi}=0{,}05$ & $12{,}4^\circ$ & $1{,}34^\circ$ & $18{,}5^\circ$ & \SI{0,28}{m} & \SI{0,9}{ms} \\
  clássica, planta ideal & $0{,}01^\circ$ & $1{,}39^\circ$ & $2{,}9^\circ$ & \SI{0,14}{m} & \SI{0,4}{ms} \\
  \bottomrule
\end{tabular}
\end{center}
A elevação residual da clássica ($\approx\SI{1,4}{\degree}$ mesmo com planta ideal) é a inclinação
que a câmera virtual esconde.

\paragraph{SITL (E3 + dois A/B, 16/11, omnisciente, 1 repetição).} Com os pesos de Sheng a clássica
mantém o alvo na imagem como a perception-aware (\SI{0,084}{} e \SI{0,070}{rad} contra \SI{0,064}{} e
\SI{0,095}{rad}; azimute 0,8--$1{,}4^\circ$), mas fica 1,2--\SI{2,1}{m} atrás: o comando satura em
\SI{2}{m/s} (= \code{max\_vel} = velocidade do alvo) e o horizonte de \SI{1}{s} copiando $\vect v_t$ não
corta caminho. $Q_{s_3}=200$ corrige a distância (0,33--\SI{0,41}{m}) mas devolve o ciclo-limite de
yaw (azimute 9--$13^\circ$, $\sigma_{\dot\psi}$ 0,43--\SI{0,73}{rad/s}, alvo fora da imagem 10--17\%), e
$R_{\dot\psi}$ até 5 não o remove. A identificação sobre esse voo dá 0,25 s + 0,16 s (confirma a malha
de yaw). Padrão mantido em $Q_{s_3}=10$. Estudos: \code{tcc/e3\_centralizacao},
\code{ibvs\_classic\_depth\_ab}, \code{ibvs\_classic\_yaw\_ab}; texto completo na monografia
(\code{latex/monografia/capitulos/ibvs\_mpc.tex}).

\section{Sintonia}
\begin{itemize}
  \item \code{yaw\_time\_constant}/\code{yaw\_delay}: identificados no SITL
        (\autoref{sec:ibvs_yaw_id}). Zerá-los volta ao integrador ideal e ao ciclo-limite.
        Em hardware real é preciso reidentificar (\code{ros2\_ws/identify\_px4\_lag.py} faz o mesmo para a
        malha de velocidade).
  \item \code{weight\_image}/\code{weight\_range}: a razão decide a troca entre centralizar e manter
        a distância. Com 400 a elevação cai para \SI{0,5}{\degree}, mas o erro de distância sobe e,
        com a câmera inclinada, o drone hesita em se aproximar (inclinar para frente desloca a imagem).
  \item \code{horizon\_N}: 12 dava elevação \SI{2}{\degree} pior com \code{weight\_vel} baixo; 16
        ($\approx\SI{2,1}{s}$) é o padrão.
  \item \code{camera\_model}: \code{body} para a câmera fixa (x500\_depth); \code{level} para o
        gimbal em \code{stabilize}. Com o gimbal em \code{track} quem centraliza é o gimbal, e esta lei
        não é necessária.
  \item \code{image\_azimuth\_deg}/\code{image\_elevation\_deg} deslocam o ponto desejado na imagem
        (por exemplo, para manter o alvo acima do centro).
\end{itemize}
